\documentclass[12pt, oneside]{book}

% metadata
\title{Interturkic Grammar}
\author{}
\date{\today}

% typography
\usepackage[skip=6pt]{caption}
\usepackage{float}
\usepackage{fontspec}
\setmainfont[Ligatures=TeX]{Cambria}
\newfontfamily{\arabicfont}[Script=Arabic, Scale=1.1]{Calibri}
\usepackage[letterpaper, top=1.5in, bottom=1.5in, left=1.5in, right=1.25in]{geometry}

% tables
\usepackage{longtable}
\usepackage{multirow}

% linguistics
\usepackage{leipzig}
\makeglossaries
\usepackage{gb4e}
\noautomath % must be called immediately after loading gb4e

\usepackage{bidi} % must be loaded last

% general settings
\setlength{\parindent}{0pt}
\setlength{\tabcolsep}{3.25pt}
\renewcommand{\arraystretch}{1.2}
\clubpenalty = 10000
\widowpenalty = 10000

% embedding of Interturkic text within English
\newcommand{\latscript}[1]{\textit{#1}}
\newcommand{\arabscript}[1]{{\arabicfont \RLE{#1}}}

% punctuation
\newcommand{\apostrophe}{\textquotesingle}
\newcommand{\zeromorpheme}{\(\emptyset\)}
\newcommand{\derivedfrom}{\(\leftarrow\)}
\newcommand{\becomes}{\(\rightarrow\)}

% tables meant to be smaller than a single page
% #1 (optional) = label
% #2 = caption
% #3 = column specification
\newenvironment{smalltable}[3][]
	{
		\begin{table}[H]
			\centering
			\caption{#2}
			#1
			\begin{tabular}{|#3|}
	}
	{
		 	\end{tabular}
		\end{table}
	}

% tables that may span multiple pages
% #1 (optional) = label
% #2 = caption
% #3 = column specification
% #4 = header
\newenvironment{bigtable}[4][]
	{
		\begin{longtable}[c]{|#3|}
			\caption{#2} #1 \\
			\hline
			#4
			\endfirsthead
			\hline
	}
	{
			\hline
		\end{longtable}
	}

\begin{document}
	\maketitle

	\frontmatter
	\setlength{\parskip}{0pt} % minimise spacing between ToC and LoT entries
	\tableofcontents
	\listoftables
	\printglosses[style=block]

	\mainmatter
	\setlength{\parskip}{6pt} % set paragraph spacing for the rest of the document
	\chapter{Preface}

\section{Intentions}

During this project, I tried to keep the following goals in mind, with the hope that they would guide me towards creating a language that would be aesthetically pleasing and internally consistent.

\begin{enumerate}
	\item Capture what I see as the basic essence of the Turkic language family as a whole, particularly with regards to features such as vowel harmony and verbal morphology;
	\item Create something that would be somewhat easy to learn and would feel familiar to the speakers of as many Turkic languages as possible;
	\item Avoid creating a mere copy of Turkish or Azerbaijani, a tendency I have seen in similar projects.
\end{enumerate}

\section{Conventions}

The conventions described in this section will be used throughout this text.
For the sake of simplicity and brevity, examples in Interturkic will be rendered only in the Latin orthography.

Slashes enclose \emph{phonemic} transcriptions, while square brackets enclose \emph{phonetic} transcriptions.
Thus, the phonemic transcription of \latscript{oquw} would be /ɔquw/, while its phonetic transcription would be [ɔquː].

Interturkic text is in italic type, often followed by the accompanying translation between single quotes (e.g.\ \latscript{Bu tilni tüşünesen mi?} `Do you understand this language?').
Glossing abbreviations are set in small caps (e.g.\ \Pl{} denotes plurality).

Dashes delineate the boundaries between morphemes, both in example text and in glosses, e.g.\ \latscript{köz-ler-imiz-den} eye-\Pl{}-\Poss{}.\Fpl{}-\Abl{} `from my eyes'.
A dash to the left of an element marks it as a verbal stem (e.g.\ \latscript{kel-} `come').
A dash to the right of an element marks it as a suffix (e.g.\ the genitive case suffix \latscript{-nIñ}).

Periods join glosses whenever they constitute the same morpheme (e.g.\ the first-person plural possessive suffix \latscript{-(I)mIz} \Poss{}.\Fpl{}), they cannot be trivially separated (e.g.\ \latscript{maña} \Fsg{}.\Dat{}), or their division is not under consideration (e.g.\ \latscript{töbeler-den} hill.\Pl{}-\Abl{} `from the hills').

Parentheses around an individual phoneme denote that its inclusion is conditional on its environment.
If the enclosed phoneme is a vowel, then it surfaces only if the suffix is joined to a consonant final stem (e.g.\ the first-person possessive suffix \latscript{-(I)m} in \latscript{qalqan-ım} `my shield and \latscript{arpa-m} `my barley').
On the other hand, if the enclosed phoneme is a consonant, then it surfaces only if the suffix is joined to a vowel final stem (e.g.\ the possessive dative case suffix \latscript{-(n)A} in \latscript{tawı-na} `to their (\Sg{}) mountain' and \latscript{közlerim-e} `to my eyes')

Square brackets around a segment indicate that its inclusion as optional (e.g.\ the accusative suffix after the third-person possessive suffix \latscript{-n[I]}).

Arrows illustrate the effect of transformative processes.
Thus, `X \becomes{} Y' would mean `X becomes Y', while `X \derivedfrom{} Y' would mean `X is derived/formed from Y'.

\end{document}
